% ===== latex-recomposed:document =====
% flm: style file flmrecomposed.b2c823f3.sty
\begin{flmFloat}{table}{NumCap}\flmCellsTable{cc}{X[-1]X[-1]}{hspan=minimal,
  cell{1}{1}={}{m, font={\flmCellsHeaderFont}},
  cell{1}{2}={}{m, font={\flmCellsHeaderFont}},
  hline{2}={1}{.4pt,solid},
  hline{2}={2}{.4pt,solid}}{\toprule
A&B\\
1&2\\
\bottomrule
}\caption{Numbers}\flmLabels{\label{ref1}}\end{flmFloat}

See \flmRef{ref1} and \flmRef[S]{ref1}.
%%
% ===== latex-recomposed:style flmrecomposed.b2c823f3.sty =====
\NeedsTeXFormat{LaTeX2e}
\ProvidesPackage{flmrecomposed.b2c823f3}
% --- packages ---
% flm.latex.pkg.caption
\RequirePackage{caption}
% flm.latex.pkg.tabularray
\RequirePackage{tabularray}
% flm.latex.pkg.xcolor
\RequirePackage{xcolor}
% flm.latex.pkg.zref-clever
\RequirePackage{zref-clever}
% --- commands and environments the library defines for the document ---
% flm.latex.def.floats
% Floats: a float of FLM's figure, table or a configured float type is
% written as \begin{flmFloat}{<type>}{<case>} <content> \caption{<caption>}
% <labels> \end{flmFloat}, <case> being NumCap (numbered, with a caption),
% NumOnly (numbered, no caption: an empty \caption{} shows the number),
% CapOnly (a caption, no number) or Bare (neither).
% \flmFloatSetUseConfig<case>{<type>}{<placement>}: how a float of the case is
%   set. By default a numbered float is set in the environment <type> at
%   <placement>, and an unnumbered one where it is written, in a center
%   environment in which \caption works as in the float <type>; a CapOnly
%   float's caption is set with no number (\caption* of the caption package).
% \flmFloatPlacementArgs<case>: the placement of a float of the case, [tbph]
%   for the numbered cases and [h] for the others.
% The content is centered. Redefine these commands to set floats differently.
\providecommand\flmFloatPlacementArgsNumCap{[tbph]}
\providecommand\flmFloatPlacementArgsNumOnly{[tbph]}
\providecommand\flmFloatPlacementArgsCapOnly{[h]}
\providecommand\flmFloatPlacementArgsBare{[h]}
\def\flm@float@infloat#1#2{\edef\flm@float@begin{\noexpand\begin{#1}#2}\def\flm@float@end{\end{#1}}}
\def\flm@float@incenter#1{\def\flm@float@begin{\begin{center}\def\@captype{#1}}\def\flm@float@end{\end{center}}}
\def\flm@float@starcaption{\let\flm@float@caption\caption\def\caption{\flm@float@caption*}}
\ProvideDocumentCommand\flmFloatSetUseConfigNumCap{m m}{\flm@float@infloat{#1}{#2}}
\ProvideDocumentCommand\flmFloatSetUseConfigNumOnly{m m}{\flm@float@infloat{#1}{#2}}
\ProvideDocumentCommand\flmFloatSetUseConfigCapOnly{m m}{\flm@float@incenter{#1}\flm@float@starcaption}
\ProvideDocumentCommand\flmFloatSetUseConfigBare{m m}{\flm@float@incenter{#1}}
\ProvideDocumentEnvironment{flmFloat}{m m}{%
  \edef\flm@float@placement{\csname flmFloatPlacementArgs#2\endcsname}%
  \csname flmFloatSetUseConfig#2\endcsname{#1}{\flm@float@placement}%
  \flm@float@begin\centering}{\flm@float@end}
% flm.latex.def.cells
% Tables of cells. A table is written
% \flmCellsTable{<natural>}{<fitted>}{<options>}{<rows>}: <natural> and
% <fitted> are two column specifications of the tblr environment of the
% tabularray package (<natural> has c for each column, <fitted> has X[-1]
% for each column),
% <options> are the table's other tblr options (the cells' spans, alignment
% and header font, and the lines), and <rows> are the rows, the cells
% separated by & and each row ended by \\.
% By default the table is set at its natural width with <natural>; where it
% is wider than \flmCellsMaxWidth, it is set again at that width with
% <fitted>. It is set between \flmCellsBeginCenter and \flmCellsEndCenter,
% except where LaTeX sets a single line, as in a cell of another table set
% at its natural width, where no paragraph can begin. A header cell is set
% in \flmCellsHeaderFont. Redefine \flmCellsTable to set tables another way,
% or one of the four other commands to change one part of it.
% A cell's custom style <name> is written \flmCellsSet{<name>} at the start
% of the cell's text: a tabularray table command that applies the tabularray
% cell keys held in \flmCellsKeys<name>. \flmCellsStyleKeys{<name>}{<keys>}
% sets those keys unless keys are already set for <name>: define
% \flmCellsKeys<name> before loading this file, or redefine it after, to
% change how a style looks. The shipped styles green, blue, yellow and red
% set the background colors flmCellsColor<name>, which you can define before
% loading this file.
% tabularray's counter library restores LaTeX's counters after the passes in
% which tabularray sets a cell only to measure it.
% The macros with @ in their names are internals, not a customization point.
\UseTblrLibrary{booktabs}
\UseTblrLibrary{counter}
\providecommand\flmCellsHeaderFont{\bfseries}
\providecommand\flmCellsBeginCenter{\begin{center}}
\providecommand\flmCellsEndCenter{\end{center}}
\providecommand\flmCellsMaxWidth{0.96\linewidth}
\ExplSyntaxOn
\long\def\flm@cells@fit#1#2#3#4{%
  \setbox0=\hbox{\begin{tblr}{colspec={#1},#3}#4\end{tblr}}%
  \dim_compare:nNnTF {\wd0} < {\flmCellsMaxWidth}
    {\leavevmode\box0}
    {\begin{tblr}{width=\flmCellsMaxWidth,colspec={#2},#3}#4\end{tblr}}}
\long\def\flm@cells@iflr#1#2{\mode_if_horizontal:TF {\mode_if_inner:TF {#1} {#2}} {#2}}
\def\flm@cells@set#1{\ifcsname flmCellsKeys#1\endcsname
  \exp_args:Nnv \tblr_set_cell:nn {} {flmCellsKeys#1}\fi}
\ExplSyntaxOff
\ProvideDocumentCommand\flmCellsTable{+m +m +m +m}{%
  \flm@cells@iflr
    {\flm@cells@fit{#1}{#2}{#3}{#4}}%
    {\flmCellsBeginCenter\flm@cells@fit{#1}{#2}{#3}{#4}\flmCellsEndCenter}}
\ProvideDocumentCommand\flmCellsStyleKeys{m m}{%
  \expandafter\providecommand\csname flmCellsKeys#1\endcsname{#2}}
\ifcsname flmCellsSet\endcsname\else
  \NewTableCommand\flmCellsSet[1]{\flm@cells@set{#1}}
  \ProvideDocumentCommand\flmCellsSet{m}{}
\fi
\providecolor{flmCellsColorgreen}{RGB}{200,255,200}
\providecolor{flmCellsColorblue}{RGB}{200,220,255}
\providecolor{flmCellsColoryellow}{RGB}{255,255,200}
\providecolor{flmCellsColorred}{RGB}{255,200,200}
\flmCellsStyleKeys{green}{bg=flmCellsColorgreen}
\flmCellsStyleKeys{blue}{bg=flmCellsColorblue}
\flmCellsStyleKeys{yellow}{bg=flmCellsColoryellow}
\flmCellsStyleKeys{red}{bg=flmCellsColorred}
% flm.latex.def.refs
% \flmRef[<options>]{<label>}: a reference that gives the name of its target
% ("Section 2"), set by \zcref of the zref-clever package.
%
% The options are the ones the document wrote, in FLM's own vocabulary. This
% definition translates each of them by its key -- a value written after the
% key is ignored, as in FLM -- and drops every key it has no translation for:
%
%   noprefix    -> noname                 the number or title alone, with no
%                                         name before it ("2", not
%                                         "Section 2")
%   S           -> S                      the first name capitalized and
%                                         unabbreviated, for a reference that
%                                         opens a sentence
%   capital     -> cap=true,abbrev=false  the names capitalized ("Lemma 1")
%   lowercase   -> cap=false,abbrev=false the names in lowercase ("lemma 1")
%   abbreviated -> abbrev=true            the abbreviated names ("Lem. 1")
%                                         where the reference type declares
%                                         them, the full names otherwise
%   anything else -> dropped              the default wording; a variant of
%                                         your own naming has no counterpart
%                                         in zref-clever, which would stop on
%                                         a key it does not know
%
% Redefine \flmRef to set references differently -- another package, another
% wording, your own variants, no link. Your definition sees the options as the
% document wrote them, in FLM's vocabulary, and maps them as it likes.
%
% The translation uses expl3, LaTeX's programming layer, which is part of the
% format and needs no package; its names read only between \ExplSyntaxOn and
% \ExplSyntaxOff. \keyval_parse:NNn reads the option list as LaTeX reads any
% key-value list: it hands a key written alone to the first macro below and a
% key with its value to the second, which drops the value. \str_case:nnF looks
% the key up in the table and appends what \zcref is to be given, or nothing.
% The macros with @ in their names are internals, not a customization point:
% redefine \flmRef, not its parts.
\ExplSyntaxOn
\def\flm@ref@add#1{\tl_put_right:Nn \flm@ref@options {#1,}}
\def\flm@ref@key#1{\str_case:nnF {#1}
  {{noprefix}{\flm@ref@add{noname}}
   {S}{\flm@ref@add{S}}
   {capital}{\flm@ref@add{cap=true,abbrev=false}}
   {lowercase}{\flm@ref@add{cap=false,abbrev=false}}
   {abbreviated}{\flm@ref@add{abbrev=true}}}
  {}}
\def\flm@ref@pair#1#2{\flm@ref@key{#1}}
\def\flm@ref@zcref#1#2{\zcref[#1]{#2}}
\ProvideDocumentCommand\flmRef{O{} m}{%
  \def\flm@ref@options{}%
  \keyval_parse:NNn \flm@ref@key \flm@ref@pair {#1}%
  \exp_args:No \flm@ref@zcref {\flm@ref@options} {#2}}
\ExplSyntaxOff
% flm.latex.def.labels
% \flmLabels{<\label commands>}: the labels of one place in the document.
% In text it runs them as written. On a line of a display-math
% environment, where amsmath keeps the last label of a line and loses the
% ones before it, it runs the first and records every further one as an alias:
% a line per alias is written to the .aux file after the last page and read on
% the next run, making the alias refer to what the first label refers to
% (for \ref and for \zcref alike). Redefine \flmLabels to change what a
% group of labels does; \def\flmLabels#1{#1} restores plain LaTeX behavior,
% amsmath's "Multiple \label's" message included.
\ProvideDocumentCommand\flmLabels{m}{%
  \ifmmode\expandafter\flm@labels@math\else\expandafter\@firstofone\fi{#1}}
% Reading the group: the first \label as written, the rest as aliases of it.
\def\flm@labels@math#1{\def\flm@lbl@first{}\flm@lbl@scan#1\flm@lbl@stop}
\def\flm@lbl@scan{\@ifnextchar\label\flm@lbl@take\flm@lbl@rest}
\def\flm@lbl@take\label#1{%
  \ifx\flm@lbl@first\@empty
    \def\flm@lbl@first{#1}%
    \label{#1}%
  \else
    \expandafter\flm@recordalias\expandafter{\flm@lbl@first}{#1}%
  \fi
  \flm@lbl@scan}
\def\flm@lbl@rest#1\flm@lbl@stop{#1}
% Collecting the aliases. A name is recorded once however often the line is
% typeset (amsmath measures a line before setting it).
\def\flm@aliaslist{}
\def\flm@recordalias#1#2{%
  \ifcsname flm@alias@#2\endcsname\else
    \expandafter\gdef\csname flm@alias@#2\endcsname{}%
    \g@addto@macro\flm@aliaslist{\flm@alias@entry{#2}{#1}}%
  \fi}
% Writing them, from the hook that runs after the last page has been shipped
% out: every \newlabel of that page has reached the .aux file by then, so the
% alias lines come after them. \AtEndDocument would run too early.
\def\flm@alias@write#1#2{\immediate\write\@auxout{\string\flm@aliaslabel{#1}{#2}}}
\AddToHook{enddocument/afterlastpage}{\let\flm@alias@entry\flm@alias@write\flm@aliaslist}
% Reading them back: the alias refers to what the original refers to. The
% assignments are global because the .aux file is read inside a group.
\def\flm@aliaslabel#1#2{%
  \ifcsname r@#2\endcsname
    \global\expandafter\let\csname r@#1\expandafter\endcsname\csname r@#2\endcsname
  \fi
  \ifcsname Z@R@#2\endcsname
    \global\expandafter\let\csname Z@R@#1\expandafter\endcsname\csname Z@R@#2\endcsname
  \fi}
% --- definitions written from the configuration of FLM features ---
% flm@float@reftype@table (the floats feature's, written from its configuration): declared unless \flm@float@reftype@table is already defined, so a definition made earlier under that name is kept
\ifcsname flm@float@reftype@table\endcsname\else\zcRefTypeSetup{table}{Name-sg={Table},name-sg={Tab.},Name-pl={Tables},name-pl={Tabs.}}\def\flm@float@reftype@table{}\fi
\endinput
